% !TEX program = pdflatex
\documentclass[11pt, a4paper]{article}

\usepackage[T1]{fontenc}
\usepackage{palatino}
\usepackage{mathrsfs}
\usepackage{amsmath, amssymb, amsthm}
\usepackage{geometry}
\usepackage{microtype}
\usepackage{setspace}
\usepackage{titlesec}
\usepackage{hyperref}
\usepackage{mdframed}
\usepackage{longtable}
\usepackage{booktabs}
\usepackage{array}
\usepackage{enumitem}
\usepackage{xcolor}

\geometry{top=1.1in, bottom=1.1in, left=1.2in, right=1.2in}
\onehalfspacing

\hypersetup{colorlinks=true, linkcolor=black, citecolor=black, urlcolor=black}

\newmdenv[linewidth=0.7pt, innertopmargin=9pt, innerbottommargin=9pt,
  innerleftmargin=12pt, innerrightmargin=12pt,
  skipabove=10pt, skipbelow=10pt]{notebox}

\newmdenv[linewidth=1.2pt, innertopmargin=10pt, innerbottommargin=10pt,
  innerleftmargin=14pt, innerrightmargin=14pt,
  skipabove=12pt, skipbelow=12pt, backgroundcolor=gray!6]{candbox}

\titleformat{\section}{\large\bfseries}{\thesection.}{0.5em}{}
\titleformat{\subsection}{\normalsize\bfseries}{\thesubsection.}{0.5em}{}
\titleformat{\subsubsection}{\normalsize\itshape}{}{}{}

\title{%
  \textbf{The Flyxion Dependency Atlas}\\[0.4em]
  \large Definitions, Theorems, Machinery, Applications,\\
  and Equivalence Candidates Across the Corpus
}
\author{Flyxion}
\date{June 2026 --- Internal Working Document. Not for distribution.}

\begin{document}
\maketitle
\thispagestyle{empty}

\begin{abstract}
This atlas inventories the theoretical corpus at four levels: primitive
definitions, major theorems and results, mathematical machinery, and
applications. Its primary purpose is to identify equivalences ---
objects that appear under different names in different frameworks but
are mathematically identical or related by a natural transformation.
Every confirmed equivalence eliminates redundant content from the
projected fifteen-volume curriculum and potentially collapses several
apparent volumes into one. The atlas is a living document updated as
new identifications are made.

\medskip
\noindent\textbf{Central thesis}: All Flyxion frameworks are constrained
gradient flows on admissibility spaces. The fifteen volumes are fifteen
instantiations of one variational principle, not fifteen independent
theories.

\medskip
\noindent\textbf{Strongest unification candidate (§5.3)}: There exists
a functor $\mathfrak{T}$ from frameworks to energy landscapes such that
$\mathfrak{T}(\text{RSVP}) = V$, $\mathfrak{T}(\text{Agency}) = F$,
$\mathfrak{T}(\text{Repair}) = d^2$, etc. If this functor exists,
Volume~XV becomes a theorem about equivalence classes of energy
functionals rather than a narrative integration.

\medskip
\noindent\textbf{Missing primitive (§1.5)}: An admissible trajectory
object $\gamma: [0,T] \to \mathcal{A}$ is absent from the current
primitive list despite being implicitly required by memory, repair,
agency, ecphory, Yarncrawler, and TARTAN. Adding it resolves multiple
apparent equivalences as instances of one path-theoretic structure.

\medskip
\noindent\textbf{Repair Theory elevation}: Repair appears across
cognitive, semantic, institutional, and physical domains. It is more
foundational than its current placement in the Application tier
suggests and should be elevated to the Core tier.
\end{abstract}

\newpage
\tableofcontents
\newpage

% ================================================================
\section{Level 1: Primitive Definitions}
% ================================================================

The following objects appear across the corpus. Each entry records the
canonical name, its mathematical type, its first formal appearance,
alternative names used elsewhere, and any suspected equivalences with
other entries.

% ----------------------------------------------------------------
\subsection{Geometric and Topological Objects}

\subsubsection{Admissibility Space / Admissibility Manifold}

\begin{longtable}{lp{10cm}}
\textbf{Symbol} & $\mathcal{A} \subseteq \mathcal{S}$ \\
\textbf{Type} & Sub-manifold (or sub-variety) of state space $\mathcal{S}$ \\
\textbf{Definition} & Subset of states satisfying all active constraints:
  $\mathcal{A} = \{x \in \mathcal{S} \mid C_i(x) \leq 0\; \forall i\}$ \\
\textbf{First appearance} & \textit{Hidden Curvature}, Chapter 2 \\
\textbf{Alternative names} & Constraint manifold (Vol.~I), admissible set
  (Institutional Attractors essay), reachable set $R(x)$ in computational
  contexts, fiber preimage $\pi^{-1}(m)$ in projection contexts,
  world-state $\Omega_t$ in Spherepop, possibility space $\mathcal{P}$
  in MEM\textbar{}8 \\
\textbf{Metric} & Fisher information metric $g_{ij} = \mathbb{E}[
  \partial_i \log p\, \partial_j \log p]$ induced by KL divergence \\
\textbf{Equivalence candidates} & \textit{See §5.1}: $\mathcal{A}$,
  $\pi^{-1}(m)$, $\Omega_t$, $\mathcal{P}$, $R(x)$ may all be instances
  of one object under different projections \\
\end{longtable}

\subsubsection{Reachability Set / Reachability Volume}

\begin{longtable}{lp{10cm}}
\textbf{Symbol} & $R(x)$; $\Omega(x) = \mathrm{Vol}(R(x))$ \\
\textbf{Type} & Subset of $\mathcal{S}$; real-valued functional \\
\textbf{Definition} & $R(x) = \{x' \in \mathcal{S} \mid \exists$
  admissible path from $x$ to $x'\}$; $\Omega(x) = \mathrm{Vol}(R(x))$ \\
\textbf{First appearance} & \textit{Computation Beyond Data}, §8 \\
\textbf{Alternative names} & Accessibility (Portuguese version),
  reachability (English version), admissible future set,
  $\mathrm{Vol}(\mathcal{A}_i)$ in social freedom formalism \\
\textbf{Note} & $R(x) \subseteq \mathcal{A}$: reachability is accessibility
  from a specific starting point, not the full admissibility space \\
\end{longtable}

\subsubsection{Projection Operator}

\begin{longtable}{lp{10cm}}
\textbf{Symbol} & $\pi: X \to M$ \\
\textbf{Type} & Smooth surjective map between manifolds \\
\textbf{Definition} & Map from full state space $X$ to observable manifold
  $M$; in general non-injective \\
\textbf{First appearance} & \textit{Hidden Curvature}, Chapter 3 \\
\textbf{Alternative names} & Compression map, visibility operator $\pi$
  (Institutional Attractors), rendering operator $\mathcal{R}$ (CLIO),
  coarse-graining map (RSVP-RG connection), realization functor $F:
  \mathcal{H} \to \mathcal{C}$ (Spherepop) \\
\textbf{Equivalence candidates} & \textit{See §5.2}: all projection-like
  maps across the corpus \\
\end{longtable}

\subsubsection{Fiber / Preimage}

\begin{longtable}{lp{10cm}}
\textbf{Symbol} & $\pi^{-1}(m) = \{x \in X \mid \pi(x) = m\}$ \\
\textbf{Type} & Subset of $X$; sub-manifold when $\pi$ is a submersion \\
\textbf{First appearance} & \textit{Hidden Curvature}, Chapter 3 \\
\textbf{Alternative names} & Preimage (CLIO), possibility space
  $\mathcal{P}_{i-1}$ (MEM\textbar{}8), event fiber (Spherepop),
  exclusion complement $\mathcal{S} \setminus \mathcal{E}$ (Spherepop) \\
\textbf{Key quantity} & $\mu(\mathcal{P}_m) \to 0$: operational fiber death \\
\end{longtable}

\subsubsection{Inadmissibility Functional}

\begin{longtable}{lp{10cm}}
\textbf{Symbol} & $\mathcal{I}(x)$; also $d(f(x), f^*)^2$ in repair context \\
\textbf{Type} & Non-negative real-valued functional on $\mathcal{S}$ \\
\textbf{Definition} & Measure of how far $x$ is from $\mathcal{A}$;
  $\mathcal{I}(x) = 0 \iff x \in \mathcal{A}$ \\
\textbf{First appearance} & Formal spine document (Vol.~I master equation) \\
\textbf{Alternative names} & Constraint violation, tension $V(x,t)$
  (Institutional Attractors), free energy $F(q,o)$ (Simulated Agency),
  merge obstruction $\mathcal{O}_\mathrm{merge}$ (Semantic Infrastructure) \\
\textbf{Note} & This object unifies ``tension,'' ``free energy,'' and
  ``constraint violation'' across the corpus. Strong equivalence candidate. \\
\end{longtable}

% ----------------------------------------------------------------
\subsection{Field Objects (RSVP)}

\subsubsection{Scalar Field}

\begin{longtable}{lp{10cm}}
\textbf{Symbol} & $\Phi(x,t)$ \\
\textbf{Type} & Scalar field on spacetime or state space \\
\textbf{Interpretation} & Structural potential / likelihood landscape /
  scalar density / salience \\
\textbf{Alternative names} & Scalar persistence $\Phi$ (MEM\textbar{}8
  memory context), scalar density (cosmological context) \\
\end{longtable}

\subsubsection{Vector Field (RSVP)}

\begin{longtable}{lp{10cm}}
\textbf{Symbol} & $\mathbf{v}(x,t)$ \\
\textbf{Type} & Vector field \\
\textbf{Interpretation} & Admissibility flow / information transport /
  lamphrodyne direction \\
\textbf{Alternative names} & Vector transport $\mathbf{v}$ (MEM\textbar{}8) \\
\end{longtable}

\subsubsection{Entropy / Accessibility Field}

\begin{longtable}{lp{10cm}}
\textbf{Symbol} & $S(x,t)$ \\
\textbf{Type} & Scalar field; in RSVP treated as a dynamical variable \\
\textbf{Interpretation} & Local accessibility density; configurational
  entropy; information compression at each scale \\
\textbf{Alternative names} & Accessibility functional $\mathcal{A}(x,t) =
  \log|\mathcal{P}(x,t)|$ (MEM\textbar{}8), opacity $\Omega$ (Institutional
  Attractors) \\
\textbf{Warning} & Three distinct uses in corpus: (a) RSVP dynamical field,
  (b) Shannon entropy $H$ in information-theoretic contexts, (c) path
  degradation measure in memory contexts. These must be carefully
  distinguished in every volume. \\
\end{longtable}

% ----------------------------------------------------------------
\subsection{Operator Objects}

\subsubsection{Pop Operator}

\begin{longtable}{lp{10cm}}
\textbf{Symbol} & $\mathsf{Pop}$ \\
\textbf{Type} & Transition operator on Spherepop configurations \\
\textbf{Effect} & Reduces local entropy $S(x_0)$ at site $x_0$ when it
  exceeds threshold $\theta_P$; records event in history \\
\textbf{Continuous analog} & Entropy reduction term in RSVP $S$-equation \\
\end{longtable}

\subsubsection{Refuse Operator}

\begin{longtable}{lp{10cm}}
\textbf{Symbol} & $\mathsf{Refuse}$ / $\mathbf{Refuse}$ \\
\textbf{Type} & Transition operator with typed annotation \\
\textbf{Effect} & Excludes trajectory $x$ with responsibility weight
  $\rho_x \neq 0$; marks event as \textsc{Refused} rather than
  \textsc{Committed} \\
\textbf{Role} & Principled exclusion preserving Hausdorff separation
  of event manifold $\mathcal{H}$ \\
\textbf{Equivalence candidates} & Constraint enforcement in constraint
  geometry; inadmissibility boundary enforcement; ``firing'' in
  threshold-based systems \\
\end{longtable}

\subsubsection{Collapse Operator}

\begin{longtable}{lp{10cm}}
\textbf{Symbol} & $\mathbf{Collapse}$ / $\mathsf{Collapse}$ \\
\textbf{Type} & Projection / realization functor $F: \mathcal{H} \to
  \mathcal{C}$ \\
\textbf{Effect} & Recovers observable state from history:
  $\mathbf{Collapse}(H_t) = T_{x_n} \circ \cdots \circ T_{x_1}(s_0)$ \\
\textbf{Alternative names} & Realization functor, MEM\textbar{}8 state
  reconstruction $s_n = \bigcap_i \mathcal{C}(e_i)$, ecphory (cognitive
  context) \\
\end{longtable}

\subsubsection{Repair Operator / Repair Morphism}

\begin{longtable}{lp{10cm}}
\textbf{Symbol} & $R_\epsilon$; formally a morphism in category of
  admissibility manifolds \\
\textbf{Type} & Map $R_\epsilon: \mathcal{S}_\mathrm{broken} \to
  \mathcal{S}$ \\
\textbf{Formal conditions (proposed)} & (a) admissibility-preserving:
  image lies in $\mathcal{A}$; (b) divergence-reducing:
  $d(f(R_\epsilon(x)), f^*) < d(f(x), f^*)$; (c) path-constrained:
  follows trajectory within current $\mathcal{A}$ \\
\textbf{Note} & This definition is the central gap in Volume~IX.
  It must be developed before that volume can be written. \\
\textbf{Alternative names} & Approximate repair, functional recovery,
  $\epsilon$-repair \\
\end{longtable}

\subsubsection{Ecphory / Memory Retrieval Operator}

\begin{longtable}{lp{10cm}}
\textbf{Symbol} & $\mathcal{E}_\mathrm{cph}$ (proposed notation) \\
\textbf{Type} & Wave propagation operator on semantic manifold \\
\textbf{Effect} & Generates retrieval wave from present state;
  wave must reach activation threshold $\theta_E$ at target state \\
\textbf{Equivalence candidates} & Repair operator restricted to cognitive
  domain; Collapse operator applied to memory history; reconstruction
  in projection theory \\
\end{longtable}

\subsubsection{Lamphrodyne Operator}

\begin{longtable}{lp{10cm}}
\textbf{Symbol} & $\mathcal{L}_\mathrm{lp}(X) =
  (-\nabla\cdot\mathbf{v},\; -\lambda\nabla\Phi,\; D_S\nabla^2 S)$ \\
\textbf{Type} & Linear relaxation operator on RSVP field triple \\
\textbf{Effect} & Separates linear lamphrodyne relaxation from nonlinear
  forcing $\mathcal{N}(X)$ in $\partial_t X = \mathcal{L}_\mathrm{lp}(X)
  + \mathcal{N}(X)$ \\
\end{longtable}

% ----------------------------------------------------------------
\subsection{Memory and History Objects}

\subsubsection{Event History}

\begin{longtable}{lp{10cm}}
\textbf{Symbol} & $H_t = (e_1, e_2, \ldots, e_n)$; also $\mu$ in MEM\textbar{}8 \\
\textbf{Type} & Finite sequence of typed events \\
\textbf{Role} & Primary object in MEM\textbar{}8 and Spherepop; state is
  derived from history, not stored directly \\
\textbf{State recovery} & $s_n = \bigcap_{i=1}^n \mathcal{C}(e_i)$ \\
\end{longtable}

\subsubsection{Memory Residue / Memory Trace}

\begin{longtable}{lp{10cm}}
\textbf{Symbol} & $M(t)$; wave memory field \\
\textbf{Type} & Continuous field of influences decaying with time \\
\textbf{Dynamics} & Damped wave equation with viscosity $\eta(x,t)$ and
  relaxation timescale $\tau(x,t)$ \\
\textbf{Alternative names} & Wave memory, collapse residue, memory state
  $M$ in formal spine Vol.~VI \\
\end{longtable}

\subsubsection{Procedural vs.\ Declarative Memory}

\begin{longtable}{lp{10cm}}
\textbf{Symbols} & $A_p(t)$, $A_d(t)$ \\
\textbf{Decay rates} & $\lambda_p \gg \lambda_d$: procedural decays faster \\
\textbf{Significance} & $A_p(t) \to 0$ while $A_d(t)$ remains:
  operational fiber death \\
\textbf{Source} & Institutional Attractors essay \\
\end{longtable}

% ----------------------------------------------------------------
\subsection{The Missing Primitive: Admissible Trajectory}

\begin{notebox}
\textbf{Finding from atlas review}: An admissible trajectory object is
absent from all formal definitions despite being implicitly required
by at least six frameworks. Its absence is the most surprising gap in
the primitive list.
\end{notebox}

\subsubsection{Admissible Trajectory}

\begin{longtable}{lp{10cm}}
\textbf{Symbol} & $\gamma: [0,T] \to \mathcal{A}$ (continuous);
  $(\gamma_0, \gamma_1, \ldots, \gamma_n) \in \mathcal{A}^{n+1}$
  (discrete) \\
\textbf{Type} & Path in admissibility space; element of trajectory
  space $\Gamma(\mathcal{A})$ \\
\textbf{Proposed definition} & A trajectory $\gamma$ is admissible if
  $\gamma(t) \in \mathcal{A}$ for all $t \in [0,T]$ and the velocity
  $\dot{\gamma}(t)$ is tangent to $\mathcal{A}$ wherever $\mathcal{A}$
  has boundary \\
\textbf{Trajectory space} & $\Gamma(\mathcal{A}) = \{\gamma: [0,T]
  \to \mathcal{A} \mid \gamma \text{ admissible}\}$ with appropriate
  topology \\
\textbf{Why it is missing} & The gradient flow master equation
  $\dot{x} = -\nabla_g \mathcal{I}(x) + \xi_t$ implicitly defines
  trajectories but never names them as a primitive. All downstream
  frameworks that use paths inherit this unnamed object. \\
\end{longtable}

\medskip
Once named, this object unifies the following as specializations:

\medskip
\begin{longtable}{lll}
\textbf{Framework object} & \textbf{Framework} & \textbf{As trajectory} \\
\hline
Memory trace $\mu = (e_1,\ldots,e_n)$ & MEM\textbar{}8 &
  Discrete trajectory in event space \\
Repair path & Repair Theory &
  Trajectory in $\mathcal{A}$ reducing $d(f(\gamma(t)), f^*)$ \\
Agency trajectory & Simulated Agency &
  Trajectory maximizing $\mathbb{E}[\Omega(\gamma(T))]$ \\
Ecphoric path & Cognitive / MEM\textbar{}8 &
  Trajectory from present state to target memory \\
Yarncrawler route & Semantic Infrastructure &
  Trajectory completing partial world-state \\
TARTAN trajectory & TARTAN &
  Trajectory preserved under decomposition \\
\end{longtable}

\medskip
\begin{notebox}
\textbf{Structural consequence}: The five-level hierarchy proposed in
the review (state spaces → admissibility → trajectory spaces →
projection structures → energy functionals → dynamics) places trajectory
spaces at Level~2, between admissibility (Level~1) and projection
(Level~3). Many frameworks differ primarily at Levels 2--4 while
sharing Levels 0--1. The trajectory object is where the hierarchy begins
to branch.
\end{notebox}



% ================================================================
\section{Level 2: Major Theorems and Results}
% ================================================================

Status codes: \textbf{Pr} = proved; \textbf{Sk} = sketch exists;
\textbf{Co} = conjecture with proof strategy; \textbf{Nu} = numerical
evidence; \textbf{As} = asserted without proof or strategy;
\textbf{Ph} = philosophical proposal; \textbf{Ga} = gap, not yet
stated precisely.

\subsection{Technical Gaps --- affect one or two volumes if failed}

\medskip
\begin{longtable}{p{5cm}lp{1.5cm}p{3.5cm}}
\toprule
\textbf{Result} & \textbf{Status} & \textbf{Home} & \textbf{Used by} \\
\midrule
Projection-Collapse Principle (weak) & Pr & I & II, III, IV, VI, VII \\
Projection-Collapse Principle (strong) & Co & I & II, III \\
Constraint-Information Equivalence & Pr & I & II, III, VIII \\
Fiber Divergence Lemma & Pr & I & III, VI \\
Reconstruction Stability ($L^2$) & Pr & I & III, VI, IX \\
Bishop-Gromov for admissibility & Pr & I & II, X \\
Persistent admissibility topology & Pr & I & VII, IX \\
Admissibility geodesics existence & Pr & I & II, X \\
Scale invariance of RSVP action & Pr & II & XIII \\
RSVP-RG equivalence & Co & II & XIII \\
Asymptotic safety as RSVP equilibrium & Co & II & XIII \\
Hausdorff preservation (Refuse) & Pr & VII & VIII \\
RSVP-Spherepop functor & Pr & VII & VIII \\
Projection-Induced Irreversibility & Pr & VII & VIII, IX \\
Operational Fiber Death & Pr & III & VI, IX \\
Reconstructability Loss & Pr & III & VI, IX \\
Approximate Repair theorem & Pr & IX & XII, XV \\
Emergent Agency theorem & Pr & IV & XV \\
Self-Concealment Principle & Pr & IV & XV \\
TARTAN structure-preservation & Pr & XI & XII \\
Yarncrawler divergence control & Pr & VIII & VIII \\
Transjective functor faithfulness & Pr & VIII & VIII \\
Coordination alignment convergence & As & X & X \\
HYDRA cognitive emergence & As & V & V, XIV \\
Spherepop confluence & Ga & VII & VII \\
HYDRA stability analysis & Ga & V & V \\
Kuramoto reduction from RSVP & Ga & X & X \\
CLIO inverse problem (ill-posedness) & Ga & III & III \\
MEM\textbar{}8 reconstruction theorem & Ga & VI & VI \\
AI Admissible Stabilization (convergence cert.) & As & III & III, XV \\
\bottomrule
\end{longtable}

\subsection{Foundational Gaps --- affect nearly every volume if failed}

\begin{notebox}
\textbf{Priority}: Attack these before writing begins. A failed
confluence theorem changes one volume. A failed Reachability Principle
changes nearly every volume. A failed Master Claim restructures
the entire project.
\end{notebox}

\medskip
\begin{longtable}{p{5cm}lp{1.5cm}p{3.8cm}}
\toprule
\textbf{Result} & \textbf{Status} & \textbf{Home} & \textbf{Risk if false} \\
\midrule
Master Claim (all systems are gradient flows) & Ph & XV & Affects all volumes \\
Reachability Principle & Sk & I & Affects IV, VI, IX, XII, XV \\
RSVP Second Law ($dS/dt \leq 0$) & As & II & Affects II, XIII \\
Lamphrodyne attractor stability & As & II & Affects II, XIII, XIV \\
Ecphoric threshold existence & As & VI & Affects IV, IX \\
Energy Functor existence ($\mathfrak{T}$) & Ph & XV & Affects XV; determines
  whether XV is synthesis or narrative \\
Invariant Discovery Principle & As & XV & Affects XV \\
Armillary Principle & Ph & XII & Affects XII \\
Optimal Opacity ($\Omega^* = \kappa/2\eta$) & Sk & IV & Affects IV (optimum proved; functional form derivation pending) \\
\bottomrule
\end{longtable}



% ================================================================
\section{Level 3: Mathematical Machinery}
% ================================================================

For each volume's content, what prior mathematics must a reader know?
This determines the prerequisites chapter of each volume.

\medskip
\begin{longtable}{p{5cm}p{8cm}}
\toprule
\textbf{Machinery} & \textbf{Where required} \\
\midrule
Riemannian geometry (manifolds, geodesics, curvature) &
  I, II, X \\
Fisher information / information geometry &
  I, II, III \\
PDEs (existence, uniqueness, regularity) &
  II, VI, XI \\
Functional analysis ($L^2$, Sobolev spaces) &
  I, II, VI \\
Dynamical systems (attractors, stability, Lyapunov functions) &
  II, IV, V, X \\
Category theory (functors, natural transformations, fibered categories) &
  I, VII, VIII \\
Sheaf theory &
  VIII \\
Optimal transport &
  I, IX, X \\
Statistical mechanics (entropy, partition functions) &
  II, XIII \\
Quantum field theory (path integrals, renormalization) &
  II, XIII, XIV \\
Operational semantics / type theory &
  VII \\
Graph theory and synchronization (Kuramoto) &
  X \\
Bayesian inference / variational methods &
  III, IV, V \\
Homotopy / persistent homology &
  I (topology section) \\
Stochastic processes / SDEs &
  I, II, IV \\
Linear algebra (SVD, Tikhonov regularization) &
  I (reconstruction section) \\
\bottomrule
\end{longtable}

\medskip
\begin{notebox}
\textbf{Finding.} Category theory and Riemannian geometry are required
by the most volumes (eight and six respectively). These two bodies of
machinery should be given full treatment in Volume~I rather than being
introduced piecemeal in later volumes. Every later volume can then
assume the reader has covered them.

The current Volume~I draft covers Riemannian geometry adequately.
The category theory chapter is the identified gap.
\end{notebox}

% ================================================================
\section{Level 4: Applications vs.\ Foundations}
% ================================================================

The following table distinguishes foundational frameworks (where new
mathematics is introduced) from application domains (where existing
mathematics is instantiated).

\medskip
\begin{longtable}{p{4cm}lp{7cm}}
\toprule
\textbf{Framework / Domain} & \textbf{Type} & \textbf{Notes} \\
\midrule
Constraint Geometry & Foundation & Volume I; source of all primitives \\
RSVP Field Theory & Foundation & Volume II; introduces PDE system \\
Projection Theory / CLIO & Foundation & Volume III; introduces inference
  pipeline \\
Spherepop Calculus & Foundation & Volume VII; introduces event language \\
Simulated Agency & Mixed & Volume IV; applies III + VI; adds agency
  formalism \\
HYDRA Architecture & Application & Volume V; applies III + V; needs
  formal coordination theory \\
MEM\textbar{}8 / Memory & Foundation & Volume VI; introduces trajectory
  memory as primary object \\
Semantic Infrastructure & Application & Volume VIII; applies VII + I \\
Repair Theory & Mixed & Volume IX; applies I + VI + VII; needs new
  repair morphism definition \\
Coordination Geometry & Application & Volume X; applies II + VIII \\
TARTAN & Mixed & Volume XI; applies I + III + VII; tiling formalism
  underdeveloped \\
Xylomorphic Computation & Application & Volume XII; applies XI; formal
  theory essentially absent \\
Cosmological Alternatives & Application & Volume XIII; applies II \\
Consciousness / Qualia & Application & Volume XIV; applies II + IV + V \\
Unified Theory & Synthesis & Volume XV; no new foundations; synthesis
  theorems only \\
\bottomrule
\end{longtable}

\medskip
\begin{notebox}
\textbf{Finding.} Five volumes are genuinely foundational (I, II, III,
VI, VII). Five are mixed (IV, IX, XI, and two others). Five are pure
applications. This means a reader who has mastered the five foundational
volumes can in principle derive the contents of the application volumes
themselves. The application volumes are valuable as guidance and
worked examples, not as sources of new mathematics.
\end{notebox}

% ================================================================
\section{Level 5: Equivalence Candidates}
% ================================================================

This is the most important section. Each entry here, if confirmed,
eliminates redundant content and potentially restructures multiple
volumes.

% ----------------------------------------------------------------
\subsection{Candidate 1: The Admissibility Cluster --- Common Parent, Not Equivalence}

\begin{candbox}
\textbf{Revised status}: The objects listed below are \emph{not} equivalent.
They occupy different categorical levels. Forcing them into a single
definition would conflate dynamical accessibility (which depends on a
starting state) with epistemic ambiguity (which depends on an observation)
and historical possibility (which depends on an event sequence). The correct
claim is weaker and more useful:

\medskip
\textbf{Claim}: $\mathcal{A}$, $R(x)$, $\pi^{-1}(m)$, $\mathcal{P}$,
$\Omega_t$ are all \emph{subobjects generated from} a common admissibility
structure. They are not the same object; they are specializations of it
under different constraints.

\medskip
\begin{tabular}{lll}
\textbf{Name} & \textbf{Framework} & \textbf{Relationship to $\mathcal{A}$} \\
\hline
Admissibility space & Vol.~I & $\mathcal{A}$ itself \\
Reachability set & Computation Beyond Data & $R(x) \subseteq \mathcal{A}$;
  path-connected component from $x$ \\
Fiber preimage & CLIO & $\pi^{-1}(m) \subseteq X$; epistemic, not dynamical \\
Possibility space & MEM\textbar{}8 & $\mathcal{P}_{i-1}$; collapsed by
  events, not by paths \\
World-state & Spherepop & $\Omega_t$; a distribution, not a set \\
\end{tabular}

\medskip
\textbf{Consequence}: Rather than a unified definition, Volume~I should
define the \emph{admissibility structure} $(\mathcal{S}, \mathcal{A},
\pi, g)$ as a tuple, and define $R(x)$, $\pi^{-1}(m)$, and $\mathcal{P}$
as derived objects with their respective constructions. This eliminates
conflation while preserving the insight that all are controlled by
admissibility.

\medskip
\textbf{Status}: Reclassified from equivalence to common parent structure.
\end{candbox}

% ----------------------------------------------------------------
\subsection{Candidate 2: The Projection Cluster}

\begin{candbox}
\textbf{Claim}: The following are all instances of the projection
operator $\pi: X \to M$.

\medskip
\begin{tabular}{lll}
\textbf{Name} & \textbf{Framework} & \textbf{Symbol} \\
\hline
Projection operator & Vol.~I & $\pi: X \to M$ \\
Rendering operator & CLIO & $\mathcal{R}$ \\
Visibility operator & Institutional Attractors & $\pi: X \to M$ \\
Coarse-graining map & RSVP-RG & (unnamed) \\
Realization functor & Spherepop & $F: \mathcal{H} \to \mathcal{C}$ \\
Compression map & Semantic Infrastructure & (various) \\
Collapse operator & MEM\textbar{}8 & $\mathbf{Collapse}$ \\
\end{tabular}

\medskip
\textbf{Key distinction}: The Spherepop realization functor and the
CLIO rendering operator both instantiate $\pi$, but they go in opposite
directions: $\pi$ maps from large to small (compression), while the
realization functor maps from history (small) to state (large).
The realization functor is closer to a \emph{section} of $\pi$ than
to $\pi$ itself. This asymmetry must be made precise before the
equivalence is confirmed.

\medskip
\textbf{Status}: Partial. The compression direction is confirmed.
The realization / section direction needs separate treatment.
\end{candbox}

% ----------------------------------------------------------------
\subsection{Candidate 3: The Tension / Free Energy Cluster}

\begin{candbox}
\textbf{Claim}: The following are all instances of the inadmissibility
functional $\mathcal{I}(x)$ --- the potential whose gradient drives
the master gradient flow in every volume.

\medskip
\begin{tabular}{lll}
\textbf{Name} & \textbf{Framework} & \textbf{Symbol} \\
\hline
Inadmissibility functional & Vol.~I & $\mathcal{I}(x)$ \\
Tension / potential & Institutional Attractors & $V(x,t)$ \\
Variational free energy & Simulated Agency & $F(q,o)$ \\
Merge obstruction & Semantic Infrastructure & $\mathcal{O}_\mathrm{merge}$ \\
Entropy drift & Semantic Infrastructure & $\mathcal{E}_\mathrm{drift}$ \\
Repair divergence & Repair Theory & $d(f(x), f^*)^2$ \\
Incoherence penalty & HYDRA & $\Omega(\mathbf{z})$ \\
\end{tabular}

\medskip
\textbf{Formal spine confirmation}: The formal spine (document 21)
confirms this cluster. Every volume's master equation has the form
$\dot{x} = -\nabla \mathcal{I}(x) + \text{terms}$. The objects in
the table are all the $\mathcal{I}$ of their respective volumes.

\medskip
\textbf{If confirmed}: this is the most important equivalence in the
corpus. It means all fifteen volumes are studying the same gradient flow
under different names and in different state spaces. Volume~XV's unified
Lagrangian $\mathcal{L}_\mathrm{Unified} = \frac{1}{2}\|\dot{\mathfrak{X}}
\|^2_\mathfrak{g} - \mathfrak{F}(\mathfrak{X})$ is already written; the
question is whether $\mathfrak{F}$ is genuinely a sum or whether some
terms are the same term appearing in different coordinates.

\medskip
\textbf{Status}: Confirmed at the structural level. Whether the
individual $\mathcal{I}$ objects are related by coordinate change or are
genuinely independent requires per-case analysis.
\end{candbox}

% ----------------------------------------------------------------
\subsection{Candidate 4: The Collapse / Reconstruction Cluster}

\begin{candbox}
\textbf{Claim}: The following are all instances of the same operation:
recovering a richer state from a compressed representation.

\medskip
\begin{tabular}{lll}
\textbf{Name} & \textbf{Framework} & \textbf{Symbol} \\
\hline
Reconstruction operator & Vol.~I & $\mathcal{R}_\epsilon$ \\
Ecphory & MEM\textbar{}8 / cognitive & $\mathcal{E}_\mathrm{cph}$ \\
Collapse (Spherepop) & Spherepop & $\mathbf{Collapse}$ \\
MEM\textbar{}8 state recovery & MEM\textbar{}8 & $s_n = \bigcap_i
  \mathcal{C}(e_i)$ \\
Approximate repair & Repair Theory & $R_\epsilon$ \\
Inverse problem solution & Vol.~III & $\hat{x}(m)$ \\
\end{tabular}

\medskip
\textbf{Key distinction}: Reconstruction in Vol.~I minimizes a
regularized inverse problem; ecphory is activation-threshold-based;
Spherepop Collapse is a functional composition; repair is a path
through admissibility space. These are not identical --- they differ in
the constraints imposed on the recovery process. But they share
the same logical form: given a compressed representation, recover
(approximately) the original.

\medskip
\textbf{Proposed unification}: A general recovery operator
$\mathcal{Q}: M \to X$ defined as a right inverse of $\pi$ under
some optimality criterion, with each specific operator corresponding
to a different criterion:
\begin{itemize}
\item Tikhonov: minimize $\|x\|^2$ subject to $\pi(x) = m$
\item Ecphory: maximize activation given viscosity-damped wave
\item Repair: stay within $\mathcal{A}$ while minimizing $d(f(x), f^*)$
\end{itemize}

\medskip
\textbf{Status}: Promising. Needs a unifying definition.
\end{candbox}

% ----------------------------------------------------------------
\subsection{Candidate 5: The Accessibility / Freedom Cluster}

\begin{candbox}
\textbf{Claim}: The following are all instances of the reachability
volume $\Omega(x) = \mathrm{Vol}(R(x))$.

\medskip
\begin{tabular}{lll}
\textbf{Name} & \textbf{Framework} & \textbf{Symbol} \\
\hline
Reachability volume & Vol.~I & $\Omega(x)$ \\
Social freedom & Institutional Attractors & $\mathcal{F} = \sum_i
  \mathrm{Vol}(\mathcal{A}_i)$ \\
Simulated reachability & Simulated Agency & $\mathcal{R}(q)$ \\
Repair bound & Repair Theory & $\eta\mathcal{R}(x)$ \\
Geometric invariance & Consciousness & $\mathcal{G}(Q)$ \\
Diagnostic volume & Institutional Attractors & $V(t)$ \\
\end{tabular}

\medskip
\textbf{Note}: The diagnostic volume $V(t)$ (number of accessible
vantage points from which a mechanism can be identified) is a
reachability volume in the space of \emph{epistemological} positions,
not physical states. This is a different type of state space but the
same mathematical structure.

\medskip
\textbf{Status}: Strong candidate. The geometric invariance
$\mathcal{G}(Q)$ in consciousness theory needs the most work to
confirm --- it may be a reachability volume in quale-state space.
\end{candbox}

% ----------------------------------------------------------------
\subsection{Candidate 6: CLIO / Agency / HYDRA / MEM|8 --- Coherent Sequence, Not One Volume}

\begin{candbox}
\textbf{Revised claim}: CLIO, Simulated Agency, HYDRA, and MEM\textbar{}8
share a common mathematical substrate but operate on different objects
and serve different pedagogical purposes. They belong in one
\emph{branch} of the curriculum, not one \emph{volume}.

\medskip
The objects genuinely differ:
\begin{itemize}
\item CLIO studies $X \xrightarrow{\pi} M$: compression and inference
\item MEM\textbar{}8 studies $(e_1,\ldots,e_n)$: sequential collapse
  and reconstruction
\item Simulated Agency studies $\arg\max_a \mathbb{E}[\Omega]$:
  future-directed optimization
\item HYDRA studies $\{f_i\}$ and coordination: distributed architecture
\end{itemize}

The analogy is probability theory, statistics, decision theory, and
control theory. They share substrate; nobody merges them into one
textbook; instead they form a coherent sequence.

\medskip
\textbf{Recommended structure}:
\begin{itemize}
\item Volume~III: Projection Theory (CLIO inference pipeline, inverse
  problems, representation entropy)
\item Volume~IV: Memory and Reconstruction (MEM\textbar{}8, ecphory,
  wave memory, trajectory residue)
\item Volume~V: Agency and Planning (Simulated Agency, free energy
  minimization, anticipated reachability)
\item Volume~VI: Modular Architectures (HYDRA coordination theory,
  stability analysis, geometric control)
\end{itemize}

The renumbering shifts the original Volumes IV--VI to a four-volume
cognitive-systems branch (new Volumes III--VI), with original Volume~VI
(MEM\textbar{}8) becoming the new Volume~IV.

\medskip
\textbf{Unifying formal object}: An agent is a tuple
$(\mathcal{A}, \pi, \mu, \mathcal{G})$ where $\mathcal{A}$ is its
admissibility space, $\pi$ its projection, $\mu$ its MEM\textbar{}8
history, and $\mathcal{G}$ its generative model. This definition
provides the common thread across all four volumes without merging them.

\medskip
\textbf{Status}: Architectural recommendation. Replaces earlier
merging proposal.
\end{candbox}

% ----------------------------------------------------------------
\subsection{Candidate 7: Repair / Yarncrawler / Semantic Infrastructure / Xylomorphic as One Scale Hierarchy}

\begin{candbox}
\textbf{Claim}: Repair theory, the Yarncrawler framework, Semantic
Infrastructure, and Xylomorphic Computation are the same
persistence-and-restoration mathematics applied at four different scales.

\medskip
\begin{tabular}{lll}
\textbf{Scale} & \textbf{Framework} & \textbf{Object being repaired} \\
\hline
Cognitive & Ecphory / MEM\textbar{}8 & Retrieval pathways \\
Semantic & Yarncrawler / Semantic Infrastructure & Knowledge structures \\
Institutional & Repair Theory & Institutional functions \\
Physical & Xylomorphic Computation & Physical infrastructure \\
\end{tabular}

\medskip
\textbf{Formal statement}: At each scale, repair is an instance of
the repair morphism $R_\epsilon$ operating on different state spaces
$\mathcal{S}$ under different constraint families $C$.

\medskip
\textbf{Consequence}: If confirmed, Volumes VIII, IX, and XII share
a common chapter on repair morphisms (citing Volume~IX for the formal
definition) and then each applies the definition to its domain.
This eliminates the gap in each volume individually and replaces it
with one central definition.

\medskip
\textbf{Status}: Strong structural case. The scale hierarchy is
real. Whether the repair morphism definition is literally the same
at each scale or merely analogous requires per-case formalization.
\end{candbox}

% ----------------------------------------------------------------
\subsection{Candidate 8: The Reachability Field as a Primitive}

\begin{candbox}
\textbf{Proposed new primitive}: Introduce a reachability field
$\mathfrak{R}(x)$ as a primitive object, with $\Omega(x) =
\mathrm{Vol}(\mathfrak{R}(x))$ as a derived quantity. Then many
objects across the corpus become instances of derivatives or
transformations of $\mathfrak{R}$.

\medskip
\begin{tabular}{ll}
\textbf{Derived quantity} & \textbf{Expression in terms of $\mathfrak{R}$} \\
\hline
Reachability volume & $\Omega(x) = \mathrm{Vol}(\mathfrak{R}(x))$ \\
Accessibility functional & $\mathcal{A}(x) = \log\Omega(x)$ \\
Inadmissibility functional & $\mathcal{I}(x) = -\log\Omega(x)$ \\
Collapse condition & $\frac{d\Omega}{dt} < 0$ \\
Repair condition & $\frac{d\Omega}{dt} > 0$ \\
Agency objective & $\arg\max_a \mathbb{E}[\Omega(x_{t+1})]$ \\
\end{tabular}

\medskip
This formulation collapses accessibility, admissibility, continuation
potential, reachability volume, and repairability into one family. The
gradient flow master equation then becomes:
\[
  \frac{dx}{dt} = \nabla_g \log\Omega(x) + \xi_t
  = -\nabla_g \mathcal{I}(x) + \xi_t,
\]
which is the Volume~I master equation rewritten in terms of
$\mathfrak{R}$. The system naturally moves toward states with
larger reachability volume.

\medskip
\textbf{If confirmed}: the reachability field $\mathfrak{R}$ replaces
several loosely related concepts with one object. The repair condition
$d\Omega/dt > 0$ becomes the formal definition of repair at all scales
(cognitive, semantic, institutional, physical), resolving the central
gap in Volume~IX without introducing new machinery.

\medskip
\textbf{Status}: Highly promising. Needs verification that
$\mathcal{I}(x) = -\log\Omega(x)$ is consistent with the existing
RSVP entropy field $S$ (which encodes similar information but is a
dynamical field, not a derived quantity). The relationship between
$\mathfrak{R}$, $S$, and $\mathcal{A}$ needs to be formalized before
Candidate~8 can be confirmed.
\end{candbox}

% ----------------------------------------------------------------
\subsection{Candidate 9: The Energy Functor $\mathfrak{T}$}

\begin{candbox}
\textbf{Claim}: There exists a functor
\[
  \mathfrak{T}: \mathbf{Frameworks} \to \mathbf{EnergyLandscapes}
\]
such that $\mathfrak{T}(\text{RSVP}) = V(x,t)$,
$\mathfrak{T}(\text{Agency}) = F(q,o)$,
$\mathfrak{T}(\text{Repair}) = d(f(x),f^*)^2$,
$\mathfrak{T}(\text{HYDRA}) = \Omega(\mathbf{z})$, etc.

\medskip
If this functor exists:
\begin{itemize}
\item The Tension/Free Energy cluster (Candidate~3) is not merely
  analogous; the individual inadmissibility functionals are related
  by morphisms in $\mathbf{EnergyLandscapes}$.
\item Volume~XV becomes a theorem about equivalence classes of energy
  functionals under $\mathfrak{T}$ rather than a narrative integration.
\item The question ``are these the same framework?'' becomes a question
  about whether $\mathfrak{T}(F_1) \cong \mathfrak{T}(F_2)$ in
  $\mathbf{EnergyLandscapes}$.
\end{itemize}

\medskip
\textbf{What $\mathbf{Frameworks}$ must be}: A category whose objects
are Flyxion frameworks and whose morphisms are admissibility-preserving
transformations between their state spaces.

\medskip
\textbf{What $\mathbf{EnergyLandscapes}$ must be}: A category whose
objects are pairs $(\mathcal{S}, \mathcal{I})$ of a state space and
an inadmissibility functional, and whose morphisms are smooth maps
$f: \mathcal{S}_1 \to \mathcal{S}_2$ satisfying
$\mathcal{I}_2(f(x)) \leq \mathcal{I}_1(x)$ (admissibility-preserving).

\medskip
\textbf{Where to spend mathematical effort}: This is the single most
productive direction for the next several months. If the functor
$\mathfrak{T}$ exists and is well-behaved, it provides the synthesis
theorem for Volume~XV and simultaneously confirms Candidate~3 and
the Master Claim.

\medskip
\textbf{Status}: Philosophical proposal. Needs formal construction.
\end{candbox}
% ================================================================
\section{The Master Claim and Its Consequences}
% ================================================================

\begin{candbox}
\textbf{Master Claim}: All Flyxion systems are constrained gradient
flows on admissibility spaces.

\medskip
\textbf{Five-level hierarchy}: Many frameworks differ primarily at
Levels 2--4 while sharing Levels 0--1. The trajectory object (§1.5)
is where the hierarchy begins to branch.

\medskip
\begin{tabular}{lll}
\textbf{Level} & \textbf{Object} & \textbf{Symbol} \\
\hline
0 & State spaces & $\mathcal{S}$ \\
1 & Admissibility structures & $\mathcal{A} \subseteq \mathcal{S}$ \\
2 & Trajectory spaces & $\Gamma(\mathcal{A})$ \\
3 & Projection structures & $\pi: X \to M$ \\
4 & Energy functionals & $\mathcal{I}: \mathcal{S} \to \mathbb{R}_{\geq 0}$ \\
5 & Dynamics & $\dot{x} = -\nabla_g \mathcal{I}(x) + \xi_t$ \\
\end{tabular}

\medskip
\textbf{Formal version}: For any Flyxion framework $\mathcal{F}$,
there exist $(\mathcal{S}_\mathcal{F}, \mathcal{A}_\mathcal{F},
\Gamma_\mathcal{F}, \pi_\mathcal{F}, g_\mathcal{F},
\mathcal{I}_\mathcal{F})$ such that:
\[
  \frac{dx}{dt} = -\nabla_{g_\mathcal{F}} \mathcal{I}_\mathcal{F}(x)
  + \xi_t + \text{source terms},
  \quad x(t) \in \Gamma_\mathcal{F}(\mathcal{A}_\mathcal{F}).
\]

\medskip
\textbf{Candidate synthesis theorem for Volume~XV}:

\begin{center}
\fbox{\parbox{0.85\textwidth}{
Every persistent system can be represented as a constrained transformation
process on a reachability structure. Formally: the category of Flyxion
frameworks embeds into $(\mathbf{AdmissibilityStructures},
\mathbf{ConstrainedGradientFlows})$.
}}
\end{center}

\medskip
\textbf{What Volume~XV must prove}: That the specific instantiations
for each framework are related by morphisms under the Energy Functor
$\mathfrak{T}$ (Candidate~9), not merely analogous.
\end{candbox}

\subsection{Repair Theory Elevation}

\begin{notebox}
\textbf{Finding}: Repair appears as a transformation class at every
scale of the corpus: ecphory (cognitive repair), Yarncrawler (semantic
repair), institutional approximate repair, and xylomorphic systems
(physical repair). The unified definition $d\Omega/dt > 0$ (Candidate~8)
makes Repair Theory a foundational result, not an application.

\medskip
\textbf{Revised dependency structure}: Volume~IX (Repair Theory) is
elevated from the Application tier to the Core tier:
\[
  \text{Vol.~I}
  \;\to\;
  \{\text{II, III, IV, VI, VII, VIII, IX}\}
\]
with Volumes V, X, XI, XII, XIII, XIV, XV as downstream consumers.
\end{notebox}



Based on the equivalence analysis, the following restructuring is
recommended before writing begins.

\subsection{Merging Candidates}

\textbf{Candidate A}: Merge Volumes IV (Simulated Agency), V (HYDRA),
and most of VI (MEM\textbar{}8) into a single volume called
\emph{Projection-Memory Agents}, with three parts:
Part~1 (MEM\textbar{}8 memory theory), Part~2 (CLIO inference
pipeline), Part~3 (HYDRA modular architecture). Simulated Agency
becomes the synthesis chapter of this merged volume. This reduces
three volumes to one.

\medskip
\textbf{Candidate B}: Merge Volume~VIII (Semantic Infrastructure)
and Volume~IX (Repair Theory) into \emph{Persistence and Repair
in Semantic Systems}. The repair morphism definition serves both.

\medskip
\textbf{Candidate C}: Reduce Volume~XII (Xylomorphic Computation)
from a full volume to two long chapters: one in the merged
Persistence and Repair volume (physical repair), one in Volume~XI
(physical tiling). Xylomorphic computation does not currently have
enough formal mathematics to sustain a full volume.

\medskip
If Candidates A, B, and C are adopted, the fifteen volumes reduce to
eleven, each with clearer thematic identity and less redundancy.

\subsection{Definite Separations}

The following volumes must remain separate regardless of merging:

\begin{itemize}
\item Volume~I (Constraint Geometry): foundational definitions only
\item Volume~II (RSVP): independent PDE system and field theory
\item Volume~VII (Spherepop): independent event calculus with distinct machinery
\item Volume~X (Coordination Geometry): Kuramoto and synchronization
  theory is sufficiently distinct
\item Volume~XIII (Cosmology): domain-specific application requiring
  separate treatment
\item Volume~XIV (Consciousness): domain-specific application
\item Volume~XV (Unified Theory): synthesis
\end{itemize}

\subsection{New Volume~0 Recommendation}

\begin{notebox}
Before any volume is written, produce a 40--60 page internal document
containing:
\begin{enumerate}
\item Canonical symbols table (every object, one notation, used
  everywhere)
\item Statement of every asserted-but-unproved result (11 items from §2)
\item Formal statement of the Master Claim with proof sketch
\item Per-equivalence-candidate: a one-page formal argument for or
  against each of the seven candidates in §5
\end{enumerate}
This document is not Volume~0 in the curriculum. It is the internal
consistency check. When it is complete, writing can begin with
confidence that the architecture is sound.
\end{notebox}

% ================================================================
\section{Symbol Consistency Requirements}
% ================================================================

The following conflicts in notation must be resolved before writing
begins. The right column gives the recommended canonical choice.

\medskip
\begin{longtable}{p{3cm}p{5cm}p{4cm}}
\toprule
\textbf{Conflict} & \textbf{Current usage} & \textbf{Recommended canonical} \\
\midrule
Entropy & $S$ (RSVP field), $H$ (Shannon), path degradation & $S$ for RSVP
  field; $H$ for Shannon throughout; $\sigma$ for path degradation \\
Admissibility & $\mathcal{A}$ (Vol.~I), $\Omega_t$ (Spherepop),
  $\mathcal{P}$ (MEM\textbar{}8) & $\mathcal{A}(\pi, \mathcal{S}, C)$
  canonical; others as derived notation \\
Reachability & $R(x)$, $\Omega(x)$, $\mathcal{F}$ (social freedom) &
  $R(x)$ for set; $\Omega(x)$ for volume; $\mathcal{F}_\mathrm{soc}$ for
  social freedom \\
Repair & $R_\epsilon$, $\mathcal{R}$, $\mathcal{R}(q)$ in agency &
  $R_\epsilon$ for repair morphism; $\mathcal{R}(q)$ for simulated
  reachability (different object) \\
Projection & $\pi$, $\mathcal{R}$ (rendering), $F$ (realization) &
  $\pi$ for compression direction; $\sigma$ for section direction;
  $F$ reserved for free energy \\
Free energy & $F(q,o)$, $\mathcal{I}$, $V(x,t)$, $\mathfrak{F}$ &
  $\mathcal{I}$ for inadmissibility in all non-thermodynamic contexts;
  $F$ for variational free energy in Bayesian contexts only \\
\bottomrule
\end{longtable}

% ================================================================
\section{State Functionals vs.\ Path Actions}
% ================================================================

\begin{notebox}
\textbf{Key finding}: Once trajectories are primitive objects (§1.5),
many frameworks naturally migrate from state-space optimization to
path-space optimization. The distinction between a state functional
$\mathcal{I}(x)$ and a path action $\mathcal{J}[\gamma]$ is not
cosmetic. It changes what the theory is fundamentally about.
\end{notebox}

\subsection{The Distinction}

A \emph{state functional} $\mathcal{I}: \mathcal{S} \to \mathbb{R}$
assigns a cost to a point. The gradient flow $\dot{x} =
-\nabla_g \mathcal{I}(x)$ describes how the system moves to reduce
cost instantaneously.

A \emph{path action} $\mathcal{J}: \Gamma(\mathcal{A}) \to \mathbb{R}$
assigns a cost to an entire trajectory:
\[
  \mathcal{J}[\gamma] = \int_0^T L(\gamma(t), \dot\gamma(t), t)\, dt.
\]
The extremal trajectory $\gamma^* = \operatorname{arg\,ext}_\gamma
\mathcal{J}[\gamma]$ is not a point-by-point descent but a globally
optimal path. The Euler-Lagrange equations for $\mathcal{J}$ govern
$\gamma^*$ and are generally different from the gradient flow of
$\mathcal{I}$ evaluated pointwise.

The two formulations coincide when $L(\gamma, \dot\gamma, t) =
\frac{1}{2}g(\dot\gamma, \dot\gamma) + \mathcal{I}(\gamma)$ and
the boundary conditions are fixed — this is Hamilton's principle.
In that case, gradient descent on $\mathcal{I}$ recovers the
Euler-Lagrange equations for $\mathcal{J}$. But in the general case,
they differ, and for the Flyxion corpus the general case is likely more
appropriate.

\subsection{Which Frameworks Naturally Use Which}

\medskip
\begin{longtable}{p{3.5cm}lp{6.5cm}}
\toprule
\textbf{Framework} & \textbf{Natural form} & \textbf{Reason} \\
\midrule
RSVP & State functional $\mathcal{I}(x,t)$ & PDE system; instantaneous
  field evolution; no memory of path \\
CLIO & State functional & Optimizes over current observation $m$;
  no trajectory structure required \\
Spherepop & Path action & Events are irreversible; the history is the
  primary object; Collapse reconstructs from history \\
MEM\textbar{}8 & Path action & Memory is a trajectory residue;
  identity is continuity of path, not current state \\
Agency & Path action & Optimizes over anticipated future trajectories,
  not current state; $\arg\max_a \mathbb{E}[\Omega(x_{t+1})]$ is a
  one-step path integral \\
Repair & Path action & Repair is a trajectory correction; the
  path must remain in $\mathcal{A}$ while reducing divergence \\
Ecphory & Path action & Retrieval wave propagates along trajectories;
  tip-of-the-tongue is a path that fails to reach threshold \\
Yarncrawler & Path action & Route completion is path optimization;
  divergence control is a path-level property \\
TARTAN & Path action & Trajectory-preserving decomposition; the
  trajectory is what is preserved, not the state \\
Coordination & Mixed & Kuramoto phases are states, but synchronization
  is a path property (convergence over time) \\
\bottomrule
\end{longtable}

\medskip
\begin{notebox}
\textbf{Finding}: RSVP and CLIO are naturally state-functional theories.
Spherepop, MEM\textbar{}8, Agency, Repair, Ecphory, Yarncrawler, and
TARTAN are naturally path-action theories. This is a genuine theoretical
distinction, not a matter of preference or notation.

The current corpus expresses all frameworks in state-functional language
because the gradient flow master equation $\dot{x} = -\nabla\mathcal{I}$
is state-based. This is an artifact of how the frameworks were developed,
not a feature of the frameworks themselves.

Translating the path-action frameworks into their natural language is
likely to produce cleaner results and expose currently hidden structure.
\end{notebox}

\subsection{The Revised Master Claim}

If the path-action frameworks are translated into their natural language,
the ultimate master claim shifts from a gradient flow statement to a
variational statement:

\medskip
\begin{center}
\fbox{\parbox{0.85\textwidth}{\centering
\textbf{Original master claim}:
$\dot{x} = -\nabla_g \mathcal{I}(x) + \xi_t$

\medskip
\textbf{Revised master claim}:
$\gamma^* = \operatorname{arg\,ext}_{\gamma \in \Gamma(\mathcal{A})}
\mathcal{J}[\gamma]$

\medskip
The gradient flow is a special case of the variational problem when
$L = \frac{1}{2}g(\dot\gamma,\dot\gamma) + \mathcal{I}(\gamma)$ and
boundary conditions are fixed.
}}
\end{center}

\medskip
The revised claim is strictly more general. It accommodates both
state-functional frameworks (where the Lagrangian is kinetic plus
potential) and path-action frameworks (where the Lagrangian encodes
path-level properties like irreversibility, memory, and coherence).

The Volume~I master equation $\dot{x} = -\nabla_g \mathcal{I}$ remains
correct as a special case and as the primary tool for Volumes II and III.
The full variational form $\gamma^* = \operatorname{arg\,ext}\mathcal{J}$
is needed for Volumes VI (MEM\textbar{}8), VII (Spherepop), IX (Repair),
and XI (TARTAN).

\subsection{New Primitive Required: Path Lagrangian}

\begin{longtable}{lp{10cm}}
\textbf{Symbol} & $L: T\mathcal{A} \times [0,T] \to \mathbb{R}$;
  $\mathcal{J}[\gamma] = \int_0^T L(\gamma, \dot\gamma, t)\,dt$ \\
\textbf{Type} & Function on the tangent bundle of admissibility space,
  possibly time-dependent \\
\textbf{Role} & Encodes the cost of being at a point $\gamma(t)$ with
  velocity $\dot\gamma(t)$ at time $t$; more general than a state potential \\
\textbf{Special cases} &
  $L = \frac{1}{2}g_{ij}\dot\gamma^i\dot\gamma^j + \mathcal{I}(\gamma)$:
  recovers gradient flow (Vol.~I, II) \newline
  $L = \mathcal{I}(\gamma) + \lambda\|\dot\gamma\|_\mathcal{H}^2$:
  viscosity-regularized path (MEM\textbar{}8, ecphory) \newline
  $L = d(f(\gamma), f^*)^2 - \eta\Omega(\gamma)$:
  repair Lagrangian (Vol.~IX) \newline
  $L = -\log p(e_t \mid H_t)$:
  Spherepop event log-likelihood (Vol.~VII) \\
\textbf{Add to Section 1} & This primitive belongs in the geometric
  objects subsection of Section~1, immediately after the admissible
  trajectory definition. \\
\end{longtable}

% ================================================================
\section{The Three-Layer Architecture}
% ================================================================

\begin{notebox}
\textbf{Most important structural finding of the atlas}: The corpus
has a three-layer architecture that is more compact than the
fifteen-volume decomposition. Once this architecture is visible,
the individual frameworks are recognizable as specializations of
three foundational theories.
\end{notebox}

\subsection{The Three Layers}

\medskip
\begin{longtable}{p{3cm}p{4cm}p{6cm}}
\toprule
\textbf{Layer} & \textbf{What it provides} & \textbf{Core object} \\
\midrule
Layer 1: Constraint Geometry &
  Admissibility spaces, metrics, and the inadmissibility functional &
  $(\mathcal{S}, \mathcal{A}, g, \mathcal{I})$ \\
Layer 2: Projection Theory &
  Observation, compression, and reconstruction &
  $\pi: X \to M$, fiber $\pi^{-1}(m)$, reconstruction operator \\
Layer 3: Trajectory Theory &
  Memory, repair, planning, and persistence &
  $\Gamma(\mathcal{A})$, $\mathcal{J}[\gamma]$, path extremals \\
\bottomrule
\end{longtable}

\subsection{How the Frameworks Distribute Across Layers}

\medskip
\begin{longtable}{p{3.5cm}lp{6.5cm}}
\toprule
\textbf{Framework} & \textbf{Primary layer} & \textbf{Note} \\
\midrule
Constraint Geometry (Vol.~I) & Layer 1 & Defines all three layers \\
RSVP (Vol.~II) & Layer 1 & Field equations on admissibility space \\
CLIO (Vol.~III) & Layer 2 & Inference as fiber navigation \\
MEM\textbar{}8 & Layer 3 & Memory as trajectory residue \\
Spherepop & Layer 3 & Event calculus as path algebra \\
Simulated Agency & Layer 2 + 3 & Projection plus anticipated trajectory \\
HYDRA & Layer 3 & Coordination of module trajectories \\
Repair Theory & Layer 3 & Trajectory correction toward function \\
Semantic Infrastructure & Layer 2 + 3 & Sheaf structure over projection
  with trajectory maintenance \\
TARTAN & Layer 3 & Trajectory-preserving decomposition \\
Coordination Geometry & Layer 1 + 3 & Coupling between trajectories
  on shared admissibility space \\
Xylomorphic Computation & Layer 3 & Physical trajectory infrastructure \\
Cosmology & Layer 1 & RSVP at cosmological scale \\
Consciousness & Layer 2 + 3 & Qualia as projection-stable trajectory
  invariants \\
Unified Theory & All & Synthesis across all three layers \\
\bottomrule
\end{longtable}

\subsection{Consequences for the Curriculum}

This three-layer picture suggests a more compact volume structure
organized around layers rather than individual frameworks:

\medskip
\begin{longtable}{p{1cm}p{4cm}p{7.5cm}}
\toprule
\textbf{Vol} & \textbf{Title} & \textbf{Content} \\
\midrule
I & Constraint Geometry &
  Layer 1: admissibility spaces, metrics, inadmissibility functional,
  gradient flows. All three layers defined here. \\
II & RSVP Field Theory &
  Layer 1 instantiated: coupled PDE system, lamphrodyne dynamics,
  cosmological and computational interpretations. \\
III & Projection and Inference &
  Layer 2: CLIO pipeline, fiber preimages, reconstruction theory,
  projection failure, representation entropy. \\
IV & Trajectory Theory &
  Layer 3 foundations: path actions, path Lagrangians, Euler-Lagrange
  on $\Gamma(\mathcal{A})$, trajectory spaces and their topology. \\
V & Memory and Reconstruction &
  Layer 3 applied: MEM\textbar{}8, wave memory, ecphory, viscosity,
  event histories as discrete trajectories. \\
VI & Agency and Planning &
  Layer 3 applied: simulated agency, anticipated reachability, free
  energy minimization, trajectory selection. \\
VII & Modular Architectures &
  Layer 3 applied: HYDRA coordination, stability analysis, distributed
  trajectory management. \\
VIII & Repair Theory &
  Layer 3 applied: repair morphisms, approximate repair, repair condition
  $d\Omega/dt > 0$, scales from cognitive to physical. \\
IX & Semantic Systems &
  Layer 2 + 3: Spherepop, Yarncrawler, Semantic Infrastructure, sheaf
  structure, trajectory completion. \\
X & Decomposition and Tiling &
  Layer 3: TARTAN, trajectory-preserving decompositions, simulation
  fidelity, annotated noise. \\
XI & Coordination Geometry &
  Layer 1 + 3: Kuramoto derivation, synchronization as trajectory
  coupling, distributed intelligence. \\
XII & Physical Systems &
  Layer 3 applied: Xylomorphic computation, repair-oriented
  infrastructure, ecological trajectories. \\
XIII & Cosmological Applications &
  Layer 1: RSVP cosmology, lamphrodyne smoothing, entropic redshift,
  comparison with standard models. \\
XIV & Consciousness and Qualia &
  Layer 2 + 3: field-theoretic consciousness, qualia as stable
  projection-trajectory invariants, neural mappings. \\
XV & Unified Theory &
  All layers: Energy Functor $\mathfrak{T}$, equivalence classes,
  synthesis theorem, open problems. \\
\bottomrule
\end{longtable}

\medskip
\begin{notebox}
\textbf{Comparison with previous structure}: This reorganization
moves from a framework-centric structure (each volume is a named
framework) to a layer-centric structure (each volume covers one
layer or one application of a layer). The number of volumes remains
fifteen, but the thematic identity of each volume is now determined
by its position in the three-layer hierarchy rather than by its
association with a named framework.

The individual framework names (CLIO, MEM\textbar{}8, HYDRA, etc.)
become chapter titles within volumes, not volume titles. This is
exactly the organizational shift that distinguishes a library of
theories from a coherent academic field.
\end{notebox}

\subsection{What Layer~4 Would Be}

The three-layer architecture raises an implicit question: is there a
Layer~4? If Layer~1 provides spaces, Layer~2 provides projections, and
Layer~3 provides trajectories, the natural next layer would provide
transformations between frameworks: the morphisms and functors that
relate one instantiation of the architecture to another. This is
precisely what the Energy Functor $\mathfrak{T}$ (Candidate~9) provides.
Layer~4, if it exists, is the content of Volume~XV.

\end{document}
